\documentclass[10pt,a4paper]{amsart}
\usepackage[margin=19mm]{geometry}
\usepackage{amsmath,amssymb,mathtools}
\usepackage[scr=rsfs,cal=euler]{mathalfa}
\usepackage{xfrac}
\usepackage[unicode,hidelinks,bookmarks=false]{hyperref}
\newcommand{\ie}{i.e.\ }
\newcommand{\cf}{cf.\ }
\newcommand{\resp}{respectively\ }
\newcommand{\Subsection}{\subsection}
\providecommand{\nobreakdash}{\nobreak\hbox{-}}
\setlength{\emergencystretch}{2em}
\multlinegap0pt
\usepackage{csquotes}
\usepackage{xcolor}
\usepackage{tensor}
\usepackage{amssymb}
\usepackage{mathtools}
\usepackage{enumerate}
\usepackage{cancel}
\usepackage{tcolorbox}
\usepackage{tikz}
\newtheorem{proposition}{Proposition}
\newtheorem{theorem}{Theorem}
\newtheorem{claim}{Claim}
\usepackage{cleveref}
\crefname{proposition}{Proposition}{Propositions}
\crefname{theorem}{Theorem}{Theorems}
\crefname{claim}{Claim}{Claims}
\DeclareMathOperator{\Ima}{Im}
\def\bR{\mathbb R}
\def\cA{\mathcal{A}}
\def\cH{\mathcal{H}}
\def\ddphi{dd^{\Phi}}
\newcommand{\vol}{\mathrm{vol}}
\title{A $dd^\Phi$-Lemma and Bott–Chern-type Cohomology for Spin(7)-Manifolds}
\author{Shubham Dwivedi, Ragini Singhal}
\date{Source: arXiv:2609.04156v1 (CC BY 4.0)}
\begin{document}
\maketitle

\noindent\emph{Source and licence.} A $dd^\Phi$-Lemma and Bott–Chern-type Cohomology for Spin(7)-Manifolds. Version arXiv:2609.04156v1 (CC BY 4.0), distributed by arXiv under the Creative Commons Attribution 4.0 International licence, http://creativecommons.org/licenses/by/4.0/, as recorded in the arXiv licence metadata for this exact version and shown on the abstract page https://arxiv.org/abs/2609.04156v1. Full licence notice: \texttt{licenses/arxiv-2609.04156-CC-BY-4.0.txt}.\par\smallskip

\noindent\textit{Editorial scope.} Counted unit: the article's theorem thm:Spin7Hodge (source0.tex lines 1453--1459). The article's complete original proof of it is delivered with it (source0.tex lines 1461--1537, one \texttt{proof} environment of 77 lines). No mathematics is written, repaired or re-derived: the statement and the proof are the article's own lines, in the article's own notation, and no retained line is substituted or re-flowed.  Retained uncounted as the source-local context the proof needs: lines 1394--1396; lines 1416--1418; lines 1447--1449.  Nothing was omitted from any retained span, so the package carries no omission record.  No retained line contains a citation, so the wrapper carries no bibliography and no bibliography entry.  No retained line contains a figure environment, an image-inclusion command or drawing code, so no image is delivered and the package declares no asset dependency.  Editorial formatting only, in addition: an AMS \texttt{amsart} preamble that restates the article's own declarations and macros used by the retained lines, namely the environments \texttt{proposition}, \texttt{theorem}, \texttt{claim} on separate counters, together with the standard packages those lines need.  The retained environments print in this document as Proposition 1, Theorem 1, Claim 1 in order of appearance, whereas the article prints the same items with the numbering of its own full document.  The complete native source of the article as distributed in the arXiv e-print for this exact version is bundled unchanged as \texttt{sources/209349/source0.tex}.
\begin{proposition}\label{prop:psymbold}
The principal symbol $\sigma(\ddphi)$ of the operator $\ddphi$ is injective.
\end{proposition}

\begin{align}\label{eq:cohospace}
K= \left(\Omega^4_1(M)\oplus \Omega^4_{35}(M) \right)\cap \ker(d).
\end{align}

\begin{align}\label{ddphiadjoint}
(\ddphi)^t: \Omega^4_1(M)\oplus \Omega^4_{35}(M)\rightarrow \Omega^0(M)\ \ \text{is}\ (\ddphi)^t(h\Phi+\eta) = -*(dd^*(h\Phi+\eta)\wedge \Phi).
\end{align}

\begin{theorem}\label{thm:Spin7Hodge}
Let $(M^8, \Phi)$ be a compact manifold with a torsion-free Spin(7)-structure $\Phi$. Then there exists an $L^2$-orthogonal decomposition of the space $K$ in \cref{eq:cohospace} as
\begin{align}\label{eq:Spin7Hodge}
K=\Ima(\ddphi) \oplus \cH^4_1\oplus \cH^4_{35},
\end{align}
where $\cH^4_1$ and $\cH^4_{35}$ denote the space of harmonic forms in the spaces $\Omega^4_1$ and $\Omega^4_{35}$ respectively.
\end{theorem}

\begin{proof}
Since the operator $\ddphi$ has an injective symbol from \Cref{prop:psymbold}, we use the general theory about such operators to find an $L^2$-orthogonal decomposition of the space $K\supseteq \Ima(\ddphi)$. Consider the Hilbert space
\begin{align*}
\cA = L^2(\Omega^4_1(M)\oplus \Omega^4_{35}(M)) \cap \ker d,
\end{align*}
where $d$ is intrepreted distributionally. As we saw in the discussions leading up to \cref{eq:cohospace} that $\Ima(\ddphi)\subset \cA$, hence we may regard
\begin{align*}
\ddphi: L^2(M) \rightarrow \cA,
\end{align*}
with domain, the Sobolev space $H^2(M)$. Thus,
\begin{align}\label{eq:hilbertspaceimage}
\cA = \Ima(\ddphi) \oplus \ker \left((dd^{\Phi})^*\mid_{\cA}\right).
\end{align}

As we are interested in the space $K$, we look at the auxiliary operator 
\begin{align*}
A&: \Omega^4_1\oplus \Omega^4_{35} \rightarrow \Omega^0 \oplus \Omega^5\\
& \alpha \mapsto \left((\ddphi)^t\alpha,\ d\alpha \right).
\end{align*}
Thus, 
\begin{align*}
\ker A =\{ \gamma \in \Omega^4_1\oplus \Omega^4_{35}\ :\ (\ddphi)^t\gamma=0, d\gamma=0\} = \ker\left( (\ddphi)^t\mid_K \right).    \end{align*}
The reason we introduce the auxiliary operator $A$ is to  guarantee smoothness of the $L^2$-forms which will be obtained by proving the injectivity of the principal symbol of the operator $A$.

We prove that the principal symbol of the operator $A$ is injective. Let $p\in M$ and suppose $\xi\in T^*_pM$ be a non-zero covector. We use the same notation $\xi$ for indentifying $\xi$ as both a vector and a covector. Using \cref{ddphiadjoint}, the principal symbol $$\sigma(A)_p(\xi)(\alpha)=\left(-*(\xi\wedge(\xi\lrcorner \alpha)\wedge \Phi), \xi\wedge \alpha \right).$$If $\alpha\in \ker(\sigma(A)_p(\xi))\subset \Omega^4_{1\oplus 35}(T^*_pM)$, then using $\xi\wedge \alpha=0$ we have
\begin{align*}
0=\xi\lrcorner(\xi \wedge \alpha) = |\xi|^2\alpha - \xi\wedge(\xi\lrcorner \alpha),
\end{align*}
which on using the fact that $*(\xi\wedge(\xi\lrcorner \alpha)\wedge \Phi)=0$ gives
\begin{align*}
0&= |\xi|^2*(\alpha \wedge \Phi),
\end{align*}
which on using the fact that for $a\Phi\in \Omega^4_1(M)$, $a\Phi\wedge \Phi=14a\vol_{\Phi}=0$ implies $a=0$, shows $\alpha \in \Omega^4_{35}(M)$. Since $\Omega^4_{-}=\Omega^4_{35}$, we have that $\alpha$ is anti-self-dual, $*\alpha=-\alpha$. Thus, again using $\xi\wedge \alpha=0$ (as $\alpha \in \ker(\sigma(A)_p(\xi)))$ implies
\begin{align*}
0=*(\xi\wedge \alpha)=\xi\lrcorner *\alpha = -\xi \lrcorner \alpha,
\end{align*}
thus proving as before,
\begin{align*}
|\xi|^2\alpha=0\ \implies\ \alpha=0.
\end{align*}
Hence the principal symbol of the auxiliary operator $A$ is injective. By standard analytic theory, $\ker A$ consists of smooth forms and as a result
\begin{align*}
\ker \left((\ddphi)^*\mid_{K} \right)=\ker \left((\ddphi)^t\mid_{K} \right).
\end{align*}

\medskip

\noindent
\begin{claim}\label{claim:Spin7Hodge}
$\ker \left((\ddphi)^t\mid_{K} \right) = \cH^4_1\oplus\cH^4_{35}$.
\end{claim}
If $\eta\in \cH^4_1\oplus\cH^4_{35}$ then in particular, $d^*\eta =0$ and hence $\eta\in \ker \left((\ddphi)^t\mid_{K} \right)$. For the converse, notice that since $K\subset \ker(d)$, hence from \cref{ddphiadjoint}
\begin{align*}
(\ddphi)^t\mid_K = -*((dd^*+d^*d)(\cdot)\wedge \Phi)=-*(\Delta(\cdot) \wedge \Phi),
\end{align*}
and thus
\begin{align*}
\ker \left((\ddphi)^t\mid_{K} \right)=\ker(\Delta(\cdot)\mid_{\Omega^4_{1\oplus 35}}\wedge \Phi)\cap K.
\end{align*}
Since $\Phi$ is torsion-free hence $\Delta(\Omega^4_{i})\subset \Omega^4_i,\ i=1,\ 35$. Let $\gamma=\gamma_1+\gamma_{35}\in \ker(\Delta(\cdot)\mid_{\Omega^4_{1\oplus 35}}\wedge \Phi)\cap K$. If $\gamma_1=a\Phi$ then $\Delta(a\Phi)\wedge \Phi=14(\Delta a)\vol_{\Phi}$ and hence $\gamma_1\in \ker(\Delta(\cdot)\mid_{\Omega^4_{1\oplus 35}}\wedge \Phi)\cap K \implies\ a=\text{constant}$ and thus
\begin{align*}
\ker(\Delta(\cdot)\mid_{\Omega^4_{1\oplus 35}}\wedge \Phi)\cap K = \bR\cdot \Phi \oplus \left(\Omega^4_{35}\cap K \right).
\end{align*}
Again, when $\gamma_{35}\in K \implies d\gamma_{35}=0$ and since $\Omega^4_{35}=\Omega^4_{-}$ hence $d^*\gamma_{35}=0$ which implies that $\gamma_{35}\in \cH^4_{35}$. Thus, $\ker(\Delta(\cdot)\mid_{\Omega^4_{1\oplus 35}}\wedge \Phi)\cap K=\cH^4_1\oplus \cH^4_{35}=\ker((\ddphi)^t\mid_K)$ which proves the claim.

\medskip

Thus, from \cref{eq:hilbertspaceimage} and \Cref{claim:Spin7Hodge}, we have
\begin{align*}
\cA = \Ima(\ddphi) \oplus \cH^4_1\oplus \cH^4_{35}.
\end{align*}
This identification is in the distributional sense and to get the equality for smooth forms, we intersect both sides of the equation with the space $K$ to get
\begin{align*}
K=\Ima(\ddphi) \oplus \cH^4_1\oplus \cH^4_{35},
\end{align*}
which proves \cref{eq:Spin7Hodge}.
\end{proof}

\end{document}
